% Exportado desde agregados; no editar cifras manualmente.
\begin{table}[!htbp]
\centering
\caption{Montos mensuales captados y referencia económica, directas 60+, 2001--2021}\label{tab:tesis-montos}
\footnotesize\setlength{\tabcolsep}{3pt}\renewcommand{\arraystretch}{1.25}
\begin{tabularx}{\linewidth}{p{.075\linewidth}p{.04\linewidth}YYYYY}
\toprule
\textbf{Ronda} & \textbf{Vía} & \textbf{N recibe} & \textbf{N nom./real} & \textbf{Mediana nominal (\$)} & \textbf{Mediana real (\$)} & \textbf{Prestación/ LPI: mediana \% (N)} \\
\midrule
2001 & K & 1\,680 & 1\,608/\allowbreak{}1\,608 & 1\,227,00 & 3\,856,54 & 80,38 (1\,170) \\
2003 & K & 1\,760 & 1\,714/\allowbreak{}1\,714 & 1\,490,00 & 4\,191,92 & 86,62 (1\,247) \\
2012 & K & 2\,819 & 2\,633/\allowbreak{}2\,633 & 2\,200,00 & 4\,233,66 & 92,84 (2\,632) \\
2015 & K & 4\,159 & 3\,965/\allowbreak{}3\,965 & 2\,200,00 & 3\,766,01 & 78,22 (3\,963) \\
2018 & K & 3\,247 & 2\,994/\allowbreak{}2\,994 & 2\,800,00 & 4\,279,56 & 95,58 (2\,994) \\
2021 & K & 3\,261 & 3\,062/\allowbreak{}3\,062 & 4\,000,00 & 5\,438,89 & 119,49 (3\,062) \\
2001 & A & 666 & 625/\allowbreak{}625 & 200,00 & 628,61 & 20,39 (81) \\
2003 & A & 895 & 821/\allowbreak{}821 & 250,00 & 703,34 & 28,87 (142) \\
2012 & A & 1\,574 & 1\,535/\allowbreak{}1\,534 & 500,00 & 962,20 & 24,09 (1\,533) \\
2015 & A & 2\,630 & 2\,582/\allowbreak{}2\,582 & 550,00 & 941,50 & 20,76 (2\,582) \\
2018 & A & 2\,764 & 2\,699/\allowbreak{}2\,699 & 580,00 & 886,48 & 20,47 (2\,699) \\
2021 & A & 4\,924 & 4\,883/\allowbreak{}4\,883 & 1\,550,00 & 2\,107,57 & 43,66 (4\,883) \\
2021 & A5 & 4\,784 & 4\,758/\allowbreak{}4\,758 & 1\,550,00 & 2\,107,57 & 43,66 (4\,758) \\
\bottomrule
\end{tabularx}
\notacuadro{K/A: suma personal observada parcial, no ingreso total. A5: evento explícito de vejez; no se añade de nuevo a A. Pesos reales de septiembre de 2026, precios del año calendario anterior a entrevista real. LPI por persona/mes del mismo periodo: la razón se calcula individualmente antes de resumir. En 2001/2003 solo localidades de al menos 100 mil habitantes tienen ámbito asignable; su N de LPI no es comparable con el total posterior. Una prestación bajo LPI no clasifica pobreza.}
\fuentecuadro{Elaboración propia, ENASEM; agregados de notebooks 06--08 y validación focalizada del informe general, secciones 10.2--10.7. Sin ponderación nacional.}
\end{table}
